% =====================================================================
%  Annexe technique : du fichier GPS brut au paramètre de Hurst
%  Description complète du pipeline avec formules et justifications
%  Compilation : xelatex annexe_H.tex
% =====================================================================
\documentclass[11pt,a4paper]{article}

\usepackage{fontspec}
\usepackage[french]{babel}
\usepackage[top=2.5cm,bottom=2.5cm,left=2.8cm,right=2.5cm]{geometry}
\usepackage{amsmath,amssymb,amsthm}
\usepackage{booktabs,array,longtable}
\usepackage{graphicx}
\usepackage{enumitem}
\usepackage{xcolor}
\usepackage{tikz}
\usetikzlibrary{arrows.meta,positioning}
\usepackage{hyperref}
\hypersetup{colorlinks=true,linkcolor=blue!50!black,urlcolor=blue!50!black,
            citecolor=blue!50!black}

\theoremstyle{definition}
\newtheorem{definition}{Définition}
\newtheorem{remarque}{Remarque}
\newtheorem{propriete}{Propriété}

\newcommand{\code}[1]{\texttt{\detokenize{#1}}}
\newcommand{\dC}[1]{\texttt{\detokenize{#1}}}  % nom de fichier
\newcommand{\R}{\mathbb{R}}
\newcommand{\E}{\mathbb{E}}

\title{\textbf{Du fichier GPS brut au paramètre de Hurst}\\[2mm]
\large Description complète du pipeline, formules et justifications}
\author{Annexe technique}
\date{Novembre 2026}

\begin{document}

\maketitle

\tableofcontents

\vspace{1cm}
\noindent\textbf{Objet.} Ce document décrit, étape par étape, le chemin qui mène
d'un fichier GPS GeoLife \code{.plt} aux paramètres de Hurst $\hat H_{\mathrm{DFA}}$,
$\hat H_{\mathrm{RS}}$ et $\hat H_{\mathrm{Whittle}}$ par segment, puis aux valeurs
« gardées » et « fortes » agrégées par utilisateur. Chaque section indique les
fichiers et fonctions du code correspondants, et justifie les formules par la
littérature. La section~\ref{sec:lois-puissance} décrit la seconde branche du
mémoire, consacrée à l'analyse des lois de puissance sur les durées de séjour et
les distances entre arrêts.

\section{Vue d'ensemble}

Le pipeline enchaîne six étapes :

\begin{enumerate}[leftmargin=1.4em]
  \item \textbf{Lecture et nettoyage} : convertir les lignes GPS brutes en une série
  temporelle valide. Fichiers : \dC{lecture.py}, \dC{nettoyage.py}, \dC{tri.py}.
  \item \textbf{Segmentation} : découper chaque trajet en segments continus. Fichier :
  \dC{segmentation.py}.
  \item \textbf{Statistiques de segment} : pas médian, distance médiane, dispersion.
  Fichier : \dC{stats_segment.py}.
  \item \textbf{Marquage des points suspects} : détecter les artefacts GPS par distance
  et par vitesse. Fichier : \dC{filtres.py}.
  \item \textbf{Sélection finale} : retenir les segments exploitables. Fichier :
  \dC{segmentation.py} (\code{filtrer_segments}, \code{filtrer_segments_stationnaires}).
  \item \textbf{Estimation et filtrage de H} : trois estimateurs, un test à mémoire
  courte, puis un filtre de fiabilité. Fichiers : \dC{hurst/*.py},
  \dC{03_estimation_H.py}, \dC{H_valides.py}.
\end{enumerate}

Les étapes 1–5 sont dans \dC{01_pretraitement.py}, l'étape 6 dans
\dC{03_estimation_H.py} et \dC{H_valides.py}.

\section{Étape 1 : lecture, nettoyage, tri}

\subsection{Format d'entrée}

Un fichier GeoLife \code{.plt} contient six lignes d'en-tête suivies d'une ligne par
point :
\[
  \underbrace{\text{latitude},}_{\phi} \;
  \underbrace{\text{longitude},}_{\lambda} \;
  0, \;
  \text{altitude (pieds)}, \;
  \text{jour julien}, \;
  \text{date}, \;
  \text{heure (GMT)}
\]
La fonction \code{lire_dossier_utilisateur} de \dC{lecture.py} convertit chaque date
et heure en secondes UNIX \emph{en UTC}. L'altitude $-777$ signale une mesure
invalide ; elle est ignorée car seules les positions horizontales entrent dans la
série.

\subsection{Nettoyage}

La fonction \code{nettoyer_points} de \dC{nettoyage.py} rejette un point si :
\begin{itemize}[leftmargin=1.4em]
  \item la latitude $\phi \notin [-90^\circ, 90^\circ]$, ou la longitude
  $\lambda \notin [-180^\circ, 180^\circ]$ ;
  \item l'horodatage est nul ou négatif.
\end{itemize}

\subsection{Tri et déduplication}

L'ordre des opérations est : nettoyage $\to$ tri par horodatage $\to$ suppression
des doublons. Parmi des points consécutifs de même horodatage, le premier est
conservé (fonction \code{supprimer_doublons} de \dC{tri.py}).

\begin{remarque}
La documentation de \code{supprimer_doublons} mentionne « même position et même
horodatage » comme critère, mais le code teste uniquement l'horodatage. Si deux
positions distinctes portent le même timestamp, la seconde est perdue. Ce choix
n'a pas d'effet mesuré sur $\hat H$.
\end{remarque}

\section{Étape 2 : segmentation temporelle}

\subsection{Règle de coupure}

Un segment est un intervalle maximal entre deux coupures. La fonction
\code{segmenter_points} de \dC{segmentation.py} coupe dès que :
\begin{equation}
  t_i - t_{i-1} > T_{\mathrm{gap}},
  \label{eq:segmentation}
\end{equation}
avec $T_{\mathrm{gap}} = \code{TIME_GAP_THRESHOLD} = 900$~s (15~min).

\subsection{Justification}

Les estimateurs de $H$ supposent une série temporellement continue (ou à
accroissements stationnaires). Un trou d'enregistrement de plusieurs minutes est un
changement de régime d'échantillonnage : la série $\{d_i\}$ contiendrait alors un
point aberrant correspondant au déplacement pendant le trou, qui n'est pas
représentatif de la dynamique intra-segment.

\begin{remarque}
Aucune interpolation n'est faite, et les segments ne sont jamais raccordés. Combler
un trou ou raccorder deux morceaux fabriquerait un saut artificiel dans la série des
distances, auquel la DFA est très sensible~\cite{chen2002}.

Le seuil de 900~s a été calibré empiriquement (voir \code{audit_grille.py}) : il
n'existe pas de « mode naturel » dans la distribution des $\Delta t$ au-delà de la
résolution native de GeoLife (1--10~s).
\end{remarque}

\section{Étape 3 : statistiques de segment}

Pour un segment de $N$ points, la fonction \code{calculer_stats_segment} de
\dC{stats_segment.py} calcule en une passe vectorisée :

\begin{itemize}[leftmargin=1.4em]
  \item les distances inter-points $d_i$ (voir §\ref{sec:distances}),
  les écarts de temps $\delta_i = t_i - t_{i-1}$,
  et les vitesses $v_i = d_i / (\delta_i / 3600)$ en km/h ;
  \item les médianes $\tilde d$ (distance) et $\tilde v$ (vitesse), plus robustes
  aux valeurs extrêmes que les moyennes ;
  \item l'écart-type $\sigma_v$ des vitesses, et le \emph{ratio d'hétérogénéité}
  \[
    R = \frac{\sigma_v}{\tilde v} ;
  \]
  \item le \emph{pas médian} $\bar{\delta} = \mathrm{med}(\delta_i)$ ;
  \item les distances médiane $\tilde d$ et maximale $d_{\max}$, la distance totale
  $D = \sum_i d_i$.
\end{itemize}

Le ratio $R$ sert uniquement à l'alerte d'hétérogénéité (voir §\ref{sec:marquage}) :
si $R > 1{,}5$, le segment est signalé comme hétérogène, mais aucun point n'est
marqué en conséquence.

\section{Étape 4 : marquage des points suspects}
\label{sec:marquage}

Le fichier \dC{filtres.py} implémente deux filtres indépendants, appliqués dans
l'ordre \code{ORDRE_FILTRES} (par défaut : distance puis vitesse).

\subsection{Filtre de distance}
\label{subsec:filtre-distance}

La fonction \code{filtrer_saut_distance} marque le point $i$ si
$d_i > d_{\max}$, où en mode relatif (V5) :
\begin{equation}
  d_{\max} = \min\bigl(30 \times \tilde d,\; 500\bigr) \quad \text{(mètres)}.
  \label{eq:filtre-distance}
\end{equation}
Le facteur $30$ est le multiplicateur de la distance médiane du segment
(\code{MAX_DIST_M = 30}), et $500$~m est le plafond absolu
(\code{MAX_DIST_M_ABSOLU = 500}).

\paragraph{Justification du mode relatif.} Un saut de 400~m dans un segment de
marche ($\tilde d \approx 3$~m) est aberrant, alors qu'un saut de 400~m dans un
segment de train ($\tilde d \approx 100$~m) ne l'est pas nécessairement. Le seuil
relatif s'adapte au contexte du segment. Le plafond absolu empêche qu'un segment à
très grand pas (avion, TGV) laisse passer un saut gigantesque sous prétexte que
$30 \times \tilde d$ serait encore plus grand.

\subsection{Filtre de vitesse}
\label{subsec:filtre-vitesse}

La fonction \code{filtrer_vitesse} implémente le Z-score modifié asymétrique
d'Iglewicz et Hoaglin~\cite{iglewicz1993}. Le Z-score standard
$z_i = (v_i - \bar v)/\sigma_v$ est sensible aux valeurs extrêmes ; le Z-score
modifié remplace la moyenne et l'écart-type par la médiane et le MAD :
\begin{equation}
  Z_i = \frac{0{,}6745\,(v_i - \tilde v)}{\mathrm{MAD}},
  \qquad
  \mathrm{MAD} = \mathrm{med}\bigl(|v_i - \tilde v|\bigr).
  \label{eq:zscore}
\end{equation}
Le facteur $0{,}6745 = \Phi^{-1}(3/4)$ est le rapport entre le MAD et l'écart-type
d'une loi normale centrée : il aligne le Z-score modifié sur l'échelle d'un z-score
standard~\cite{leys2013}.

\paragraph{Asymétrie.} Le Z-score n'est calculé que pour $v_i > \tilde v$ :
\[
  Z_i = \begin{cases}
    \dfrac{0{,}6745\,(v_i - \tilde v)}{\mathrm{MAD}} & \text{si } v_i > \tilde v, \\[2mm]
    0 & \text{sinon.}
  \end{cases}
\]
Une vitesse basse (arrêt, ralentissement) n'est jamais un artefact : elle est
légitime dans une trajectoire. Seules les vitesses \emph{supérieures} à la médiane
peuvent être candidates.

\paragraph{Seuil candidat.} Un point est candidat si
\[
  Z_i > \code{MAD_Z_SEUIL} = 3{,}5
  \quad \text{ou} \quad
  v_i > v_{\max},
\]
où $v_{\max}$ est le seuil absolu (V5 : $\min(200, \max(5, 20\tilde v))$~km/h).

\paragraph{Classification des candidats.} Un candidat isolé de variation brutale
$\Delta v = |v_i - v_{i-1}| > \max(20, 5\,\mathrm{MAD}_v)$ est un \emph{artefact}
(glitch GPS ponctuel). Un cluster de candidats consécutifs de longueur
$\geq \max(5, \mathrm{round}(0{,}01\,N))$ est un \emph{régime} (changement de mode
de transport). Dans les deux cas, seul le premier point du cluster est marqué
(\code{type_suspect_vitesse = "artefact"} ou \code{"regime"}).

\paragraph{Alerte d'hétérogénéité.} Si $R = \sigma_v/\tilde v > 1{,}5$, une alerte
globale \code{alerte_heterogene = True} est portée par le segment entier. Aucun
point n'est marqué en conséquence : c'est une note pour l'interprétation.

\begin{remarque}
Les points \code{regime} ne sont jamais des points de coupure dans
\code{couper_aux_points_marques}. Seuls les points \code{artefact} et
\code{suspect_dist} déclenchent une coupure. La variante
$\hat H_{\mathrm{morceaux}}$ (section §\ref{sec:morceaux}) utilise cette coupure
\emph{sans raccorder} les morceaux, ce qui évite le défaut de la V4 (raccordement
qui fabrique un saut).
\end{remarque}

\section{Étape 5 : sélection finale des segments}

Deux filtres retiennent les segments exploitables. La fonction
\code{filtrer_segments} impose :
\[
  N \geq \code{MIN_POINTS} = 300.
\]
La fonction \code{filtrer_segments_stationnaires} impose :
\[
  D = \sum_{i=1}^{N-1} d_i \geq \delta_{\min}(\text{segment}),
\]
où en mode relatif (V5) :
\begin{equation}
  \delta_{\min}(\text{segment}) = \max\left(
    \code{DISTANCE_MIN},\;
    \frac{\code{MIN_POINTS} \times \bar{\delta} \times V_{\mathrm{plancher}}}{3600}
  \right).
  \label{eq:distance-min}
\end{equation}
avec \code{DISTANCE_MIN} $= 0{,}1$~km et $V_{\mathrm{plancher}} = 1$~km/h
(vitesse de « marche lente »).

\paragraph{Justification.} Un segment à pas médian $5$~s et un segment à pas médian
$60$~s n'ont pas le même sens physique pour le même « 0,1~km ». La formule
\eqref{eq:distance-min} exige qu'un segment ait parcouru au moins la distance qu'un
marcheur lent aurait parcourue dans le même nombre de pas. Sur GeoLife
($\bar{\delta} \approx 7$~s), le seuil effectif passe de $0{,}1$~km à
$0{,}583$~km : environ $5{,}8$ fois plus strict.

\begin{remarque}
Le nom \code{filtrer_segments_stationnaires} est trompeur : ce n'est pas un test
de stationnarité statistique, c'est un filtre de distance minimale. Le nom vient
d'une ancienne version où l'on assimilait « immobile » à « stationnaire ». Le vrai
test de stationnarité est celui de la section §\ref{sec:statio} (mais il n'est pas
implémenté dans le code actuel).
\end{remarque}

\section{Construction de la série des distances}
\label{sec:distances}

Pour un segment de $N$ points $\{(p_1, \dots, p_N)\}$ avec
$p_i = (\phi_i, \lambda_i, t_i)$, la série transmise aux estimateurs est
$\{d_1, d_2, \dots, d_{N-1}\}$ où $d_i$ est la distance de haversine entre
$p_i$ et $p_{i+1}$ :
\begin{equation}
  a_i = \sin^2\!\Bigl(\frac{\phi_{i+1} - \phi_i}{2}\Bigr)
        + \cos\phi_i \cos\phi_{i+1} \sin^2\!\Bigl(\frac{\lambda_{i+1} - \lambda_i}{2}\Bigr),
  \quad
  c_i = 2 \arctan_2\bigl(\sqrt{a_i}, \sqrt{1-a_i}\bigr),
  \quad
  d_i = R \cdot c_i,
  \label{eq:haversine}
\end{equation}
avec $R = 6\,371$~km (rayon terrestre moyen). La fonction
\code{haversine_vectorise} de \dC{geodesie.py} implémente \eqref{eq:haversine}
en numpy ; $d_i$ est exprimé en mètres.

\begin{remarque}[Série $d_i$ : intensité, pas direction]
La série $d_i$ mesure l'intensité du déplacement à chaque pas de temps, pas la
direction. Un $H$ élevé traduit la persistance de périodes de mouvement intense
ou d'immobilité, pas celle de la trajectoire dans l'espace. La série mélange aussi
vitesse et intervalle de temps, puisque $d_i = v_i \cdot \delta_i$. L'option
\code{--serie vitesse} utilise $v_i$ pour mesurer si les conclusions dépendent de
ce choix.
\end{remarque}

% =====================================================================
\section{Étape 6a : estimation de H par segment}
\label{sec:estimation}
% =====================================================================

La fonction \code{estimer_segment} de \dC{estimation_segment.py} applique
successivement les trois estimateurs et le test nul à la série $x = \{d_1, \dots,
d_{N-1}\}$ (ou $x = \{v_i\}$ si \code{--serie vitesse}). Nous décrivons chaque
estimateur.

\subsection{Préliminaire : profil cumulé (DFA)}

Soit $x_1, \dots, x_n$ la série d'entrée ($n = N-1$), et $\bar x$ sa moyenne.
Le \emph{profil cumulé} est défini par :
\begin{equation}
  Y(k) = \sum_{i=1}^{k}(x_i - \bar x), \qquad k = 1, \dots, n.
  \label{eq:profil}
\end{equation}
On a $Y(n) = 0$ par construction. $Y(k)$ représente l'écart cumulé à la moyenne,
c'est-à-dire la « marche » de la série.

\subsection{DFA : Detrended Fluctuation Analysis}

\subsubsection{Échelles}

On choisit $n_s = 20$ échelles $s_1 < s_2 < \dots < s_{n_s}$ réparties
logarithmiquement entre $s_{\min} = 4$ et $s_{\max} = \lfloor n/4 \rfloor$ :
\begin{equation}
  s_j = \left\lfloor \exp\!\left(
    \log s_{\min} + \frac{j-1}{n_s - 1}(\log s_{\max} - \log s_{\min})
  \right) \right\rfloor, \qquad j = 1, \dots, n_s.
  \label{eq:echelles-dfa}
\end{equation}
Les doublons introduits par $\lfloor \cdot \rfloor$ sont supprimés. La fonction
\code{echelles_dfa} de \dC{dfa.py} implémente \eqref{eq:echelles-dfa}.

\paragraph{Justification des bornes.} $s_{\min} = 4$ est le minimum pour un
détrendage linéaire (deux paramètres à estimer sur au moins deux points par
paramètre, plus une marge)~\cite{peng1994}. $s_{\max} = n/4$ garantit au moins
$4$ fenêtres disjointes à la plus grande échelle, ce qui réduit la variance de
$F(s_{\max})$.

\subsubsection{Fluctuations}

Pour chaque échelle $s_j$, le profil $Y$ est découpé en
$M_j = \lfloor n/s_j \rfloor$ fenêtres disjointes de taille $s_j$ :
\[
  \left[Y((\nu-1)s_j + 1), \dots, Y(\nu s_j)\right],
  \qquad \nu = 1, \dots, M_j.
\]
Dans chaque fenêtre, on ajuste par moindres carrés un polynôme d'ordre 1
(\code{DFA_ORDRE = 1}) :
\[
  Y((\nu-1)s_j + k) \approx a_\nu\, k + b_\nu,
  \qquad k = 1, \dots, s_j.
\]
La \emph{variance locale} des résidus est :
\[
  F^2(s_j, \nu) = \frac{1}{s_j}\sum_{k=1}^{s_j}
  \bigl[Y((\nu-1)s_j + k) - (a_\nu k + b_\nu)\bigr]^2.
\]
La \emph{fluctuation} à l'échelle $s_j$ est la moyenne quadratique sur les
fenêtres :
\begin{equation}
  F(s_j) = \sqrt{\frac{1}{M_j} \sum_{\nu=1}^{M_j} F^2(s_j, \nu)}.
  \label{eq:F-s}
\end{equation}

\subsubsection{Régression log-log}

Si $F(s) \sim s^{\alpha}$, on estime $\alpha$ par régression linéaire des
$\log F(s_j)$ sur les $\log s_j$ :
\begin{equation}
  \log F(s_j) = \alpha \log s_j + c + \varepsilon_j.
  \label{eq:regression-dfa}
\end{equation}
L'estimateur des moindres carrés est :
\[
  \hat \alpha = \frac{n_s \sum_j u_j v_j - \sum_j u_j \sum_j v_j}
                     {n_s \sum_j u_j^2 - (\sum_j u_j)^2},
  \qquad u_j = \log s_j, \quad v_j = \log F(s_j).
\]

\paragraph{Interprétation.} Pour un \emph{bruit gaussien fractionnaire} (fGn,
stationnaire), $\alpha = H$. Pour un \emph{mouvement brownien fractionnaire}
(fBm, non stationnaire), $\alpha = H + 1$~\cite{peng1994}. Une valeur
$\hat\alpha > 1$ sur une série qu'on suppose stationnaire signale donc une
non-stationnarité (régime, tendance) : c'est le drapeau \code{dfa_sup_1}.

\subsubsection{Diagnostics complémentaires}

La fonction \code{dfa_complet} calcule également :
\begin{itemize}[leftmargin=1.4em]
  \item les pentes $\hat\alpha_{\mathrm{bas}}$ et $\hat\alpha_{\mathrm{haut}}$
  sur les moitiés basse et haute des échelles ;
  \item l'écart $\Delta = \hat\alpha_{\mathrm{haut}} - \hat\alpha_{\mathrm{bas}}$
  (\code{ecart_pentes}) ;
  \item les pentes locales sur fenêtres glissantes de 5 échelles
  (\code{pentes_locales}).
\end{itemize}
L'écart $\Delta$ est un diagnostic de \emph{crossover} : $\Delta \approx 0$ indique
une loi d'échelle unique (une seule pente sur toute la plage), $|\Delta|$ grand
indique deux régimes d'échelle distincts (par exemple, un régime à petite échelle
et un autre à grande échelle). Cette décomposition répond directement à l'objection
de Maraun, Rust et Timmer~\cite{maraun2004} : une pente globale ajustée sur toute
la courbe ne prouve pas une loi d'échelle.

\subsection{R/S : Rescaled Range}

\subsubsection{Statistique sur un bloc}

Pour un bloc de taille $m$ extrait de la série, soit $\bar x_m$ sa moyenne :
\[
  \bar x_m = \frac{1}{m}\sum_{i=1}^{m} x_i.
\]
L'écart cumulé au sein du bloc est :
\[
  W(k) = \sum_{i=1}^{k}(x_i - \bar x_m), \qquad k = 1, \dots, m.
\]
L'\emph{amplitude} et l'\emph{écart-type} sont :
\[
  R(m) = \max_{1 \leq k \leq m} W(k) - \min_{1 \leq k \leq m} W(k),
  \qquad
  S(m) = \sqrt{\frac{1}{m}\sum_{i=1}^{m}(x_i - \bar x_m)^2}.
\]
La statistique $R/S$ du bloc est $R(m)/S(m)$. La fonction \code{rs_bloc} de
\dC{rs.py} implémente cette définition.

\paragraph{Convention.} Le diviseur de l'écart-type est $m$ (variance biaisée),
conformément à Hurst~\cite{hurst1951} et à Mandelbrot et Wallis~\cite{mandelbrot1969}.
Certaines variantes utilisent $m-1$ ; le choix doit être cohérent avec la formule
d'Anis-Lloyd.

\subsubsection{Estimation de H}

La fonction \code{rs_complet} calcule $R/S$ pour $n_m = 20$ tailles de bloc
$m_1 < \dots < m_{n_m}$, réparties logarithmiquement entre $m_{\min} = 10$ et
$m_{\max} = \lfloor n/2 \rfloor$. Pour chaque $m_j$, on moyenne $R/S$ sur tous les
blocs disjoints de taille $m_j$ du segment, en ignorant les blocs d'écart-type nul :
\[
  \overline{R/S}(m_j) = \frac{1}{|\mathcal{B}_j^{>0}|}
    \sum_{B \in \mathcal{B}_j^{>0}} \frac{R(m_j, B)}{S(m_j, B)},
\]
où $\mathcal{B}_j^{>0}$ est l'ensemble des blocs de taille $m_j$ d'écart-type non
nul.

Si $\E[R/S] \sim c\,m^H$, on estime $H$ par régression linéaire des
$\log \overline{R/S}(m_j)$ sur $\log m_j$ :
\[
  \log \overline{R/S}(m_j) = H \log m_j + c' + \varepsilon_j.
\]
C'est l'estimation \emph{brute} (\code{h_rs_brut}).

\subsubsection{Correction Anis-Lloyd}

Le $R/S$ brut est biaisé pour les petits échantillons~\cite{anis1976}. Pour un
bruit blanc gaussien i.i.d.\ de taille $m$, l'espérance exacte de $R/S$ est :
\begin{equation}
  \E_{\mathrm{iid}}\!\left[\frac{R}{S}\right](m) =
  \frac{\Gamma\bigl((m-1)/2\bigr)}{\sqrt{\pi}\,\Gamma(m/2)}
  \sum_{i=1}^{m-1} \sqrt{\frac{m-i}{i}}.
  \label{eq:anis-lloyd}
\end{equation}
La fonction \code{esperance_rs} de \dC{rs.py} implémente \eqref{eq:anis-lloyd}
exactement pour $m \leq 340$, et utilise l'approximation asymptotique
$\Gamma((m-1)/2)/\Gamma(m/2) \approx \sqrt{2/m}$ au-delà.

La correction Anis-Lloyd estime $H$ par :
\begin{equation}
  \hat H_{\mathrm{corrigé}} = \frac{1}{2} + \text{pente}\left[
    \log \overline{R/S}(m_j) - \log \E_{\mathrm{iid}}\!\left[\frac{R}{S}\right](m_j)
  \right] \text{ sur } \log m_j.
  \label{eq:rs-corrige}
\end{equation}
C'est l'estimation corrigée (\code{h_rs_corrige}).

\paragraph{Limite mesurée.} La correction d'Anis-Lloyd est conçue pour retirer le
biais à $H = 1/2$. Sur des séries \emph{très persistantes} ($H \geq 0.7$), elle
\emph{sur-corrige} : le biais observé est de $-0.149$ à $H_0 = 0.9$, $N = 300$ (voir
tableau~\ref{tab:exactitude} du mémoire). C'est pourquoi le R/S corrigé n'est pas
utilisé comme référence indépendante pour valider un $H$ élevé.

\subsection{Whittle local}
\label{sec:whittle}

\subsubsection{Périodogramme}

Soit $x_1, \dots, x_n$ la série. Le \emph{périodogramme} est :
\begin{equation}
  I_n(\lambda_j) = \frac{1}{2\pi n} \left| \sum_{k=1}^{n} x_k e^{i k \lambda_j} \right|^2,
  \qquad \lambda_j = \frac{2\pi j}{n}, \quad j = 1, \dots, \lfloor n/2 \rfloor.
  \label{eq:periodogramme}
\end{equation}
La fonction \code{periodogramme_rapide} de \dC{whittle.py} implémente
\eqref{eq:periodogramme} via la FFT de numpy (\code{np.fft.rfft}).

\subsubsection{Critère de Whittle local}

L'estimateur de Whittle \emph{local} de Robinson~\cite{robinson1995} n'utilise que
les basses fréquences, là où le spectre d'un processus à longue portée se comporte
comme $f(\lambda) \sim c\,|\lambda|^{1-2H}$. On retient les $m$ premières
fréquences, avec :
\begin{equation}
  m = \min\left(\bigl\lfloor n^{0{,}65} \bigr\rfloor,\; \bigl\lfloor n/4 \bigr\rfloor \right).
  \label{eq:m-whittle}
\end{equation}
Le critère de Whittle à minimiser sur $H$ est :
\begin{equation}
  R(H) = \log\!\left(\frac{1}{m}\sum_{j=1}^{m} \lambda_j^{2H-1} I_n(\lambda_j)\right)
         - (2H-1)\,\frac{1}{m}\sum_{j=1}^{m} \log \lambda_j.
  \label{eq:whittle}
\end{equation}

\paragraph{Justification.} La relation $f(\lambda) \sim c\,|\lambda|^{1-2H}$ aux
basses fréquences est la définition spectrale de la dépendance à longue
portée~\cite{robinson1995}. Le critère \eqref{eq:whittle} est une version réduite
du critère de Whittle : la constante $c$ a été éliminée analytiquement en
concentrant la vraisemblance sur $H$.

\subsubsection{Minimisation}

On minimise $R(H)$ sur une grille régulière de $K = 1000$ points :
\[
  H_k = \frac{k}{K+1}, \qquad k = 1, \dots, K.
\]
L'estimation est $\hat H = H_{k^*}$ où $k^* = \arg\min_k R(H_k)$.

\paragraph{Indicateur de saturation.} Quand le minimum est atteint à une borne de
la grille ($k^* = 1$ ou $k^* = K$), cela signifie que la vraie solution serait
$H \leq 0$ ou $H \geq 1$ : la valeur n'est pas une estimation fiable. Le drapeau
\code{whittle_a_la_borne = True} est alors levé, et le segment est exclu du
filtre de garde.

\paragraph{Limite.} Le Whittle local n'est asymptotiquement normal que pour
$H < 7/8$~\cite{phillips2004}. Il est \emph{inconsistant} pour $H \geq 1$. Ce n'est
pas le Whittle exact de Shimotsu et Phillips~\cite{shimotsu2005}, qui n'est pas
implémenté.

\subsection{Test nul à mémoire courte}
\label{sec:test-nul}

Pour chaque segment, on teste si la pente DFA observée $\hat h_{\mathrm{obs}}$ peut
être expliquée par un modèle à mémoire courte de même longueur et de même
distribution. La fonction \code{evaluer_depassement_dfa} de
\dC{validation/modeles_nuls.py} implémente le test.

\subsubsection{Gaussianisation}

La série $x$ est transformée en scores normaux par ses rangs :
\begin{equation}
  z_i = \Phi^{-1}\!\left(\frac{\mathrm{rang}(x_i) - 0{,}5}{n}\right),
  \label{eq:gaussianisation}
\end{equation}
où $\Phi^{-1}$ est la fonction quantile de la loi normale standard, et le rang
$\mathrm{rang}(x_i) \in \{1, \dots, n\}$ est déterminé par ordre croissant. Cette
transformation donne une série approximativement gaussienne dont les corrélations
de rang sont préservées.

\subsubsection{Ajustement AR(p)}

On ajuste un modèle autorégressif d'ordre $p$ aux données gaussianisées $z$ :
\begin{equation}
  z_t = \phi_1 z_{t-1} + \phi_2 z_{t-2} + \dots + \phi_p z_{t-p} + \varepsilon_t,
  \qquad \varepsilon_t \sim \mathcal{N}(0, \sigma^2).
  \label{eq:ar}
\end{equation}
Les paramètres $(\phi, \sigma^2)$ sont estimés par les équations de Yule-Walker :
\begin{equation}
  \Gamma_p \phi = \gamma_{1:p}, \qquad \sigma^2 = \gamma_0 - \phi^\top \gamma_{1:p},
  \label{eq:yule-walker}
\end{equation}
où $\Gamma_p$ est la matrice de Toeplitz des autocovariances $\gamma(0), \dots,
\gamma(p-1)$ et $\gamma_{1:p} = (\gamma(1), \dots, \gamma(p))^\top$. La fonction
\code{ajuster_ar} utilise \code{scipy.linalg.solve_toeplitz} pour résoudre
\eqref{eq:yule-walker}.

L'ordre $p$ est choisi par minimisation du critère d'information d'Akaike :
\[
  \mathrm{AIC}(p) = n \log \sigma_p^2 + 2p, \qquad p \in \{0, 1, \dots, \code{NUL_ORDRE_AR_MAX}\}.
\]

\subsubsection{Simulation et échange de rangs}

On simule $B = 99$ séries AR($p$) indépendantes de même longueur $n$, avec les
paramètres $(\phi, \sigma^2)$ estimés. Chaque série simulée est ensuite transformée
pour lui redonner exactement la distribution empirique de $x$, par échange de
rangs :
\[
  s^{(b)}_i = x_{(\mathrm{rang}(z^{(b)}_i))},
\]
où $x_{(1)} \leq x_{(2)} \leq \dots \leq x_{(n)}$ est la série triée de $x$, et le
rang est celui de la simulation gaussienne.

\paragraph{Justification.} Le test vise à isoler la mémoire longue \emph{au-delà}
de ce qu'un AR($p$) peut expliquer. L'échange de rangs garantit que la marginale
de $x$ (queues lourdes, asymétrie, valeurs positives) est préservée dans les
surrogats : la seule différence entre $x$ et $s^{(b)}$ est la structure de
dépendance temporelle.

\subsubsection{Statistique et p-valeur}

On calcule la pente DFA $\hat h^{(b)}$ sur chaque surrogate $s^{(b)}$. La p-valeur
unilatérale (H$_0$ = mémoire courte) est :
\begin{equation}
  p = \frac{1 + \#\{b : \hat h^{(b)} \geq \hat h_{\mathrm{obs}}\}}{B + 1}.
  \label{eq:p-valeur-nul}
\end{equation}
Le \emph{seuil} est le quantile à $95\%$ des pentes simulées :
\[
  \mathrm{seuil}_{95} = Q_{0{,}95}\bigl(\{\hat h^{(b)}\}\bigr).
\]
L'\emph{excès} est :
\[
  \mathrm{excès} = \hat h_{\mathrm{obs}} - \mathrm{seuil}_{95}.
\]
Un excès positif indique un dépassement significatif au seuil $\alpha = 0{,}05$.

\begin{remarque}
Ce test est \emph{unilatéral} : il ne détecte que la persistance
($\hat h_{\mathrm{obs}} > \hat h^{(b)}$), pas l'anti-persistance. Ce n'est pas un
test de mémoire longue : un dépassement indique qu'un AR($p$) ne suffit pas, mais
une alternance de régimes le produit aussi~\cite{diebold2001}.
\end{remarque}

\subsection{Variante par morceaux}
\label{sec:morceaux}

La fonction \code{h_dfa_par_morceaux} de \dC{estimation_segment.py} calcule
$\hat H_{\mathrm{DFA}}$ après avoir coupé le segment aux points
\code{artefact} et \code{suspect_dist}, \emph{sans} raccorder les morceaux. Chaque
morceau de $\geq 50$ points est analysé séparément, puis les $\hat H$ sont combinés
par moyenne pondérée :
\[
  \hat H_{\mathrm{morceaux}} = \frac{\sum_{m \in \text{morceaux}} n_m \hat H_m}
                                    {\sum_{m} n_m},
\]
où $n_m$ est le nombre de points du morceau $m$.

\paragraph{Justification.} La V4 calculait un « H sans suspects » en \emph{retirant}
les points suspects puis en recalculant les distances entre voisins restants.
Cela fabriquait un saut artificiel entre les deux voisins de part et d'autre du
point retiré, saut auquel la DFA est sensible~\cite{chen2002}. La variante par
morceaux coupe sans raccorder.

% =====================================================================
\section{Étape 6b : filtre de fiabilité}
\label{sec:filtre}
% =====================================================================

La fonction \code{traiter_utilisateur} de \dC{H_valides.py} lit les résultats
\code{resultats_H/<user>_H.csv} produits par \dC{03_estimation_H.py} et
applique le filtre d'inclusion.

\subsection{Les quatre conditions}

Un segment est \emph{gardé} si et seulement si les quatre conditions suivantes sont
satisfaites :
\begin{enumerate}[leftmargin=1.4em]
  \item \textbf{Calculabilité.} Les trois méthodes ont produit une valeur :
  $\hat H_{\mathrm{DFA}}$, $\hat H_{\mathrm{RS}}$, $\hat H_{\mathrm{Whittle}}$
  ne sont ni \code{NaN} ni vides.
  \item \textbf{Domaine.} Chacune est dans $]0, 1[$ :
  \[
    0 < \hat H_{\mathrm{DFA}} < 1, \quad
    0 < \hat H_{\mathrm{RS}} < 1, \quad
    0 < \hat H_{\mathrm{Whittle}} < 1.
  \]
  \item \textbf{Whittle non saturé.} \code{whittle_a_la_borne = False}
  (section §\ref{sec:whittle}).
  \item \textbf{Accord inter-méthodes.} La divergence
  \[
    \mathrm{div} = \max(\hat H_{\mathrm{DFA}}, \hat H_{\mathrm{RS}},
                         \hat H_{\mathrm{Whittle}})
                  - \min(\hat H_{\mathrm{DFA}}, \hat H_{\mathrm{RS}},
                         \hat H_{\mathrm{Whittle}})
  \]
  vérifie $\mathrm{div} \leq 0{,}15$ (\code{SEUIL_DIVERGENCE}).
\end{enumerate}

\paragraph{Condition sur \code{whittle_a_la_borne}.} Le booléen doit être
strictement \code{False}. Une valeur vide (\code{None}) signifie que Whittle n'a
pas été calculé : le segment est exclu (pas de comparaison possible). Cette
distinction est codée explicitement par la fonction \code{_parse_bool}.

\subsection{Critique de la circularité}

La condition (4) introduit une \emph{circularité} : la divergence est calculée à
partir des $\hat H$ eux-mêmes, et le seuil $0{,}15$ conditionne l'inclusion du
segment. La convergence observée sur les segments gardés n'est donc pas une
confirmation indépendante de la qualité des estimateurs : elle est la
\emph{conséquence} du critère d'inclusion.

Sur l'échantillon de GeoLife, l'écart moyen
$|\hat H_{\mathrm{DFA}} - \hat H_{\mathrm{RS}}|$ passe de $0{,}074$ (tous les
segments) à $0{,}023$ (segments gardés) : cette réduction est attendue par
construction.

\subsection{Scénario \emph{forts}}

Parmi les segments gardés, un sous-ensemble est qualifié de \emph{fort} si les
trois méthodes donnent simultanément $\hat H > 0{,}5$ :
\[
  \hat H_{\mathrm{DFA}} > 0{,}5 \quad \text{et} \quad
  \hat H_{\mathrm{RS}} > 0{,}5 \quad \text{et} \quad
  \hat H_{\mathrm{Whittle}} > 0{,}5.
\]
La colonne \code{fort_0p5} dans \code{<user>_H_gardes.csv} encode ce booléen.

\paragraph{Interprétation.} Sur le corpus GeoLife, la quasi-totalité des segments
gardés est également forte ($n_{\mathrm{forts}} / n_{\mathrm{gardés}} \approx 0{,}97$).
Cela signifie que le filtre de garde sélectionne majoritairement des segments
persistants. Ce n'est pas une preuve de mémoire longue (le test nul n'est pas
utilisé ici), mais c'est cohérent avec l'hypothèse d'une persistance locale.

\subsection{Agrégation par utilisateur}

Pour chaque utilisateur, la fonction \code{traiter_utilisateur} calcule deux
moyennes sur les segments gardés (arithmétique et pondérée par $n_{\text{gardés}}$),
pour les trois méthodes et les trois scénarios (tous, gardés, forts) :
\[
  \bar h_{\mathrm{user}} = \frac{1}{|\mathcal{G}|}\sum_{s \in \mathcal{G}} \hat h_s,
  \qquad
  \bar h^{\mathrm{pond}}_{\mathrm{user}} =
  \frac{\sum_{s \in \mathcal{G}} n_s \hat h_s}{\sum_{s \in \mathcal{G}} n_s},
\]
où $\mathcal{G}$ est l'ensemble des segments gardés de l'utilisateur, et $n_s$ le
nombre de points du segment $s$. Le script \dC{H_valides.py} présente la moyenne
pondérée comme \emph{la} métrique de référence (les gros enregistreurs pèsent plus).

% =====================================================================
\section{Seconde branche : analyse des lois de puissance}
\label{sec:lois-puissance}
% =====================================================================

\paragraph{Objet.} Les sections précédentes décrivent la branche qui mesure une
propriété de \emph{dépendance temporelle} (l'exposant de Hurst) sur la série des
distances inter-points $d_i$. La branche décrite ici mesure une propriété de
\emph{distribution} sur les mêmes trajectoires : les durées d'arrêt et les
distances entre arrêts suivent-elles une loi de puissance ? Les deux branches sont
indépendantes dans leur traitement mais partagent la même détection d'arrêts.

\subsection{Variables analysées}
\label{sec:arrets-branche-lp}

Les deux variables sont construites par le module \dC{pretraitement/arrets.py} à
partir des CSV produits par \dC{01_pretraitement.py} :

\begin{itemize}[leftmargin=1.4em]
  \item \textbf{temps de séjour} $T$ : durée d'un arrêt, détecté par fenêtre
  glissante de Li et al.\ (rayon $R_{\mathrm{arr}}$, durée minimale
  $T_{\min}$) ;
  \item \textbf{distance entre arrêts} $D$ : distance haversine entre les centroïdes
  de deux arrêts consécutifs.
\end{itemize}

Le module distingue quatre variantes transmises au module d'analyse :
\code{sejour} (arrêts internes non censurés, choix principal),
\code{distance} (paires dans le même segment), \code{sejour_avec_censurees}
(incluant les pauses entre segments) et \code{distance_a_cheval} (incluant les
paires séparées par un trou d'enregistrement).

\begin{remarque}[Censure]
Suivant Kivelä et Porter~\cite{kivela2015}, les arrêts dont le début ou la fin
tombent hors de la fenêtre d'observation sont \emph{étiquetés} (colonne
\code{censuree}), pas supprimés. La variante \code{sejour} exclut les censurés
(car les longs séjours sont sous-représentés, ce qui biaise la queue), et
\code{sejour_avec_censurees} les inclut en sensibilité.
\end{remarque}

% ---------------------------------------------------------------------
\subsection{Étape 1 : ajustement de Clauset (loi de puissance pure)}
\label{sec:clauset}

On suppose que la queue de la distribution (au-dessus d'un seuil $x_{\min}$)
suit une loi de puissance continue :
\begin{equation}
  p(x) = \frac{\alpha - 1}{x_{\min}} \left(\frac{x}{x_{\min}}\right)^{-\alpha},
  \qquad x \geq x_{\min}, \quad \alpha > 1.
  \label{eq:loi-puissance}
\end{equation}
La fonction de survie est $P(X > x) = (x/x_{\min})^{1-\alpha}$.

\subsubsection{Maximum de vraisemblance à $x_{\min}$ fixé}

Pour un échantillon $x_1, \dots, x_n \geq x_{\min}$, la log-vraisemblance de
\eqref{eq:loi-puissance} est :
\[
  \ell(\alpha) = n \log(\alpha - 1) - n \log x_{\min}
                 - \alpha \sum_{i=1}^{n} \log(x_i / x_{\min}).
\]
L'annulation de la dérivée par rapport à $\alpha$ donne l'estimateur :
\begin{equation}
  \hat\alpha = 1 + n \left[\sum_{i=1}^{n} \log(x_i / x_{\min})\right]^{-1},
  \qquad
  \hat\sigma_\alpha = \frac{\hat\alpha - 1}{\sqrt{n}}.
  \label{eq:mle-alpha}
\end{equation}
La fonction \code{exposant_vraisemblance_max} de \dC{lois_de_puissance/clauset.py}
implémente \eqref{eq:mle-alpha}.

\subsubsection{Choix de $x_{\min}$ par distance KS}

Le seuil $x_{\min}$ est choisi en minimisant la distance de
Kolmogorov-Smirnov entre la queue observée et la loi ajustée :
\begin{equation}
  \hat{x}_{\min} = \arg\min_{x_{\min}} \sup_{x \geq x_{\min}}
  \left| \hat F_{\text{emp}}(x) - F(x \mid \hat\alpha(x_{\min}), x_{\min}) \right|,
  \label{eq:choix-xmin}
\end{equation}
où $\hat F_{\text{emp}}$ est la fonction de répartition empirique de la queue, et
$F$ la fonction de répartition théorique de \eqref{eq:loi-puissance}.

\begin{remarque}[Grille de candidats]
Le package \code{powerlaw} de Alstott et al.\ essaie \emph{toutes} les valeurs
distinctes de l'échantillon. Le code implémenté limite à
\code{LP_NB_CANDIDATS_XMIN} $= 200$ seuils équi-répartis, avec la contrainte
que la queue fasse au moins \code{LP_NB_MIN_QUEUE} $= 50$ points. Ce compromis
réduit le coût du bootstrap (qui re-teste $x_{\min}$ à chaque simulation) et
change très peu la valeur retenue.
\end{remarque}

\subsubsection{Test d'adéquation par bootstrap semi-paramétrique}

Un bon ajustement KS ne prouve pas que la loi est correcte : il faut tester
l'adéquation. La méthode de Clauset, Shalizi et Newman~\cite{clauset2009} simule
$B$ échantillons qui suivent la loi ajustée et refait tout l'ajustement (y compris
le choix de $x_{\min}$), puis compare la distribution des distances KS simulées à
celle observée :
\begin{equation}
  p_{\text{boot}} = \frac{1}{B} \sum_{b=1}^{B}
  \mathbb{1}\bigl[\mathrm{KS}_b \geq \mathrm{KS}_{\text{obs}}\bigr].
  \label{eq:p-boot-clauset}
\end{equation}
La règle est : $p_{\text{boot}} \leq 0{,}1$ $\Rightarrow$ loi de puissance
rejetée ; $p_{\text{boot}} > 0{,}1$ $\Rightarrow$ non rejetée (plausible, pas
prouvée). La fonction \code{p_valeur_adequation} de \dC{clauset.py} implémente
\eqref{eq:p-boot-clauset} avec $B = 1000$.

% ---------------------------------------------------------------------
\subsection{Étape 2 : quatre lois alternatives}
\label{sec:alternatives}

La loi de puissance pure est mise en concurrence avec quatre alternatives, toutes
ajustées sur la \emph{même} queue ($x \geq x_{\min}$) et renormalisées pour ne
décrire que cette queue. Le module \dC{lois_de_puissance/alternatives.py}
implémente chacune.

\subsubsection{Exponentielle décalée (1 paramètre)}

\[
  p(x) = \lambda \exp\bigl(-\lambda(x - x_{\min})\bigr), \qquad x \geq x_{\min}.
\]
L'estimateur du maximum de vraisemblance est explicite :
$\hat\lambda = 1/\overline{(x - x_{\min})}$.

\subsubsection{Log-normale conditionnée (2 paramètres)}

\[
  p(x \mid x \geq x_{\min}) = \frac{f(x)}{P(X \geq x_{\min})},
\]
où $f$ est la densité log-normale de paramètres $(\mu, \sigma)$ :
\[
  f(x) = \frac{1}{x \sigma \sqrt{2\pi}} \exp\left(-\frac{(\log x - \mu)^2}{2\sigma^2}\right).
\]
Le dénominateur est calculé par la fonction de survie de la loi normale
$\Phi((\log x_{\min} - \mu)/\sigma)$, en logarithme pour rester précis loin dans
la queue (\code{stats.norm.logsf} de scipy). L'ajustement se fait numériquement
par Nelder-Mead sur $(\mu, \ln\sigma)$.

\subsubsection{Weibull conditionnée (2 paramètres)}

\[
  p(x \mid x \geq x_{\min}) = f(x) \cdot \exp\bigl((x_{\min}/\lambda)^k\bigr),
\]
où $f$ est la densité Weibull de forme $k$ et d'échelle $\lambda$ :
\[
  f(x) = \frac{k}{\lambda} \left(\frac{x}{\lambda}\right)^{k-1}
         \exp\left(-\left(\frac{x}{\lambda}\right)^k\right).
\]
L'ajustement se fait par Nelder-Mead sur $(\ln k, \ln\lambda)$. La forme $k = 1$
redonne l'exponentielle ; $k < 1$ donne une queue plus lourde.

\subsubsection{Loi de puissance tronquée (2 paramètres)}
\label{sec:tronquee}

C'est l'alternative principale : elle contient la loi de puissance pure comme cas
limite et permet de tester formellement si la queue est \emph{coupée} au-delà
d'une certaine échelle :
\begin{equation}
  p(x) = \frac{1}{Z}\, x^{-\alpha} \exp(-\lambda x), \qquad x \geq x_{\min},
  \label{eq:tronquee}
\end{equation}
avec $Z = \int_{x_{\min}}^{\infty} t^{-\alpha} \exp(-\lambda t)\, dt$.

\paragraph{Calcul de $Z$ par changement de variable.} Le calcul direct de $Z$ est
numériquement instable quand $\lambda \to 0$ (l'intégrale devient proche de celle
de la loi de puissance pure, dominée par $t^{-\alpha}$). On effectue deux
changements de variables :
\begin{enumerate}[leftmargin=1.4em]
  \item $x = x_{\min} + s/\lambda$, d'où $x^{-\alpha} = x_{\min}^{-\alpha}
  (1 + s/(\lambda x_{\min}))^{-\alpha}$ et
  $\exp(-\lambda x) = e^{-\kappa} \exp(-s)$ avec $\kappa = \lambda x_{\min}$ ;
  \item $s = \kappa e^w$, d'où $ds = \kappa e^w\, dw$.
\end{enumerate}
On obtient :
\begin{equation}
  Z = \frac{x_{\min}^{-\alpha} e^{-\kappa}}{\lambda}
      \int_{-\infty}^{+\infty}
      \exp\Bigl(-\alpha \log(1 + e^w) - \kappa e^w + \log\kappa + w\Bigr)\, dw.
  \label{eq:Z-tronquee}
\end{equation}
L'intégrande en $w$ est une bosse lisse de largeur d'ordre 1, quelle que soit la
valeur de $\lambda$ : c'est ce qui rend l'intégration numérique stable même quand
$\lambda$ est très petit. La fonction
\code{_log_normalisation_puissance_tronquee} de \dC{alternatives.py}
implémente \eqref{eq:Z-tronquee}.

\paragraph{Ajustement.} Le maximum de vraisemblance de $(\alpha, \lambda)$ se
calcule par Nelder-Mead en partant de $\alpha_0$ (exposant de la loi pure) et
d'une coupure faible $\lambda_0 = 1/(10 \max x_i)$. L'optimisation ne peut donc
qu'améliorer le modèle emboîté : si la coupure n'apporte rien, $\lambda$ tend
vers $0$ et le résultat rejoint la loi de puissance pure.

% ---------------------------------------------------------------------
\subsection{Étape 3 : comparaison des vraisemblances}
\label{sec:comparaison}

Les cinq modèles sont comparés par le \emph{rapport de log-vraisemblance} :
\begin{equation}
  R = \sum_{i=1}^{n} \bigl[\log p_{\text{PL}}(x_i) - \log p_{\text{alt}}(x_i)\bigr].
  \label{eq:rapport-vraisemblance}
\end{equation}
$R > 0$ favorise la loi de puissance, $R < 0$ favorise l'alternative. La
significativité de $R$ dépend du type de modèle :

\subsubsection{Test de Vuong (alternatives non emboîtées)}

Pour les trois alternatives non emboîtées (exponentielle, log-normale, Weibull),
on utilise le test de Vuong~\cite{vuong1989} :
\begin{equation}
  p_{\text{Vuong}} = \mathrm{erfc}\!\left(
    \frac{|R|}{\sqrt{2n}\,\sigma_{\Delta}}
  \right),
  \qquad
  \sigma_{\Delta} = \mathrm{std}\bigl(\log p_{\text{PL}}(x_i) - \log p_{\text{alt}}(x_i)\bigr).
  \label{eq:vuong}
\end{equation}
Si $p_{\text{Vuong}} \leq 0{,}1$, la différence de vraisemblance est significative ;
sinon la comparaison est indécidable.

\subsubsection{Test du $\chi^2$ (alternative emboîtée : tronquée)}

La loi tronquée contient la loi de puissance pure comme cas limite
($\lambda = 0$). Le test approprié est alors un test du $\chi^2$ à un degré de
liberté sur $2|R|$ :
\begin{equation}
  p_{\text{emboîté}} = \chi^2_1\bigl(2|R|\bigr).
  \label{eq:chi2}
\end{equation}

\begin{remarque}[Paramètre au bord]
$\lambda = 0$ est \emph{au bord} de l'espace des paramètres ($\lambda \geq 0$),
ce qui rend \eqref{eq:chi2} un peu conservateur (Self \& Liang, 1987). Une
correction exacte multiplierait la p-valeur par $0{,}5$. On retient la version
conservatrice, cohérente avec la littérature Clauset.
\end{remarque}

\subsubsection{Critère d'Akaike corrigé (AICc)}

En complément des rapports de vraisemblance, on calcule les \emph{poids d'Akaike}
corrigés pour petits échantillons, qui pénalisent le nombre de paramètres et
somment à 1 sur les cinq modèles :
\begin{equation}
  \mathrm{AICc}_m = 2k_m - 2\ln L_m + \frac{2k_m(k_m + 1)}{n - k_m - 1},
  \label{eq:aicc}
\end{equation}
où $k_m$ est le nombre de paramètres du modèle $m$ et $n$ la taille de la queue.
Le poids est ensuite :
\[
  w_m = \frac{\exp(-\Delta_m/2)}{\sum_{m'} \exp(-\Delta_{m'}/2)},
  \qquad \Delta_m = \mathrm{AICc}_m - \min_{m'} \mathrm{AICc}_{m'}.
\]
La fonction \code{poids_aicc} de \dC{alternatives.py} implémente \eqref{eq:aicc}.

% ---------------------------------------------------------------------
\subsection{Étape 4 : adéquation de la loi tronquée}
\label{sec:adequation-tronquee}

Le critère AICc dit quel modèle est le meilleur ; il ne dit pas si ce modèle est
un \emph{bon} ajustement. La même distinction qu'à l'étape 1 (test KS) s'applique
donc à la loi tronquée gagnante.

L'étape 4 est un test bootstrap analogue à \eqref{eq:p-boot-clauset}, mais
spécifique à la tronquée :
\begin{enumerate}[leftmargin=1.4em]
  \item calculer la distance KS entre les données et la tronquée ajustée
  (avec $x_{\min}$ fixé à sa valeur observée) ;
  \item simuler $B$ échantillons de la tronquée (méthode par rejet :
  tirer selon la loi de puissance pure et accepter avec probabilité
  $\exp(-\lambda(x - x_{\min}))$) ;
  \item réajuster chaque simulation (avec $x_{\min}$ toujours fixé) et
  calculer sa distance KS ;
  \item $p = \#\{\mathrm{KS}_b \geq \mathrm{KS}_{\text{obs}}\} / B$.
\end{enumerate}
La fonction \code{p_valeur_adequation_tronquee} de \dC{alternatives.py}
implémente ce test avec $B = 200$.

\begin{remarque}[Puissance limitée]
Le test est \emph{beaucoup moins puissant} que celui de Clauset, pour deux
raisons : $x_{\min}$ est fixé (pas de sensibilité au mauvais choix de seuil), et
seuls $(\alpha, \lambda)$ sont réajustés sur chaque simulation. Un
$p > 0{,}1$ ne prouve donc pas que la tronquée est adéquate ; il indique
seulement que les données ne s'en écartent pas de manière détectable à cette
résolution. Un $p \leq 0{,}1$ est en revanche un rejet robuste (le test est
permissif, donc s'il rejette, c'est que le désaccord est net).
\end{remarque}

% ---------------------------------------------------------------------
\subsection{Étape 5 : lien avec l'exposant de Hurst (processus ON/OFF)}
\label{sec:onoff-annexe}

Si l'activité alterne entre phases ON (mouvement) et OFF (arrêt) dont les durées
ont une queue de survie en $t^{-a}$ avec $1 < a < 2$, on attend une dépendance
de longue portée avec~\cite{taqqu1997} :
\begin{equation}
  H = \frac{3 - a}{2} = \frac{4 - \alpha}{2},
  \qquad 2 < \alpha < 3.
  \label{eq:onoff-annexe}
\end{equation}
La condition $2 < \alpha < 3$ vient de $a = \alpha - 1$ et $1 < a < 2$ (exposant
de survie). L'exposant $\alpha$ de la densité est celui ajusté à l'étape 1.

Le code calcule automatiquement :
\begin{itemize}[leftmargin=1.4em]
  \item \code{dans_domaine_on_off} : booléen $2 < \alpha < 3$ ;
  \item \code{h_on_off_predit} : $(4 - \alpha)/2$ si dans le domaine,
  \code{NaN} sinon.
\end{itemize}
Ces deux colonnes servent à vérifier \emph{avant} toute comparaison à $H$ que la
relation est applicable : si $\alpha$ est hors domaine, la prédiction ne
s'applique pas et doit être rapportée telle quelle.

\begin{remarque}
Le résultat de Taqqu, Willinger et Sherman~\cite{taqqu1997} porte sur la
\emph{superposition} d'un grand nombre de sources ON/OFF ; pour la trajectoire
d'un seul individu, l'argument relève de la théorie du renouvellement. La
relation \eqref{eq:onoff-annexe} est donc une hypothèse de travail, pas un
théorème.
\end{remarque}

% ---------------------------------------------------------------------
\subsection{Tableau récapitulatif des constantes LP}
\label{sec:constantes-lp}

\begin{table}[h]
\centering
\caption{Constantes de la branche lois de puissance (\code{config.py}, V5)}
\label{tab:constantes-lp}
\small
\begin{tabular}{@{}p{5cm}p{1.8cm}p{7cm}@{}}
\toprule
Constante & Valeur & Rôle \\
\midrule
\code{LP_SEUIL_P} & 0{,}1 & seuil de rejet (Clauset et al., §5) \\
\code{LP_NB_BOOTSTRAP} & 1000 & simulations bootstrap pour $p_{\text{boot}}$ \\
\code{LP_NB_MIN_QUEUE} & 50 & taille minimale de la queue \\
\code{LP_NB_CANDIDATS_XMIN} & 200 & seuils $x_{\min}$ essayés \\
\code{LP_NB_MIN_PAR_UTILISATEUR} & 100 & observations min par user \\
\code{ARRET_RAYON_M} & 200 & rayon d'arrêt (m) \\
\code{ARRET_DUREE_MIN_S} & 1200 & durée minimale d'arrêt (s) \\
\code{ARRET_TROU_MAX_S} & 900 & trou qui coupe la série (s) \\
\code{ARRET_RAYON_PAUSE_M} & 200 & rayon d'une pause entre segments (m) \\
\code{ARRET_DISTANCE_MIN_M} & 100 & distance minimale entre deux arrêts (m) \\
\bottomrule
\end{tabular}
\end{table}

% ---------------------------------------------------------------------
\subsection{Chiffres de référence sur GeoLife}
\label{sec:resultats-lp-annexe}

L'analyse a été exécutée sur les 137 utilisateurs de GeoLife ayant au moins un
arrêt détecté (jeu de seuils 200~m~/~20~min), pour les quatre variables. Les
quatre résultats agrégés (population regroupée) sont reportés au chapitre 9 du
mémoire. On retiendra :
\begin{itemize}[leftmargin=1.4em]
  \item la loi de puissance pure est rejetée pour 3 variables sur 4
  (\code{distance}, \code{sejour_avec_censurees}, \code{distance_a_cheval}) ;
  \item la loi de puissance tronquée obtient le meilleur poids AICc pour
  \code{sejour} (0{,}97) et \code{distance_a_cheval} (0{,}51) ;
  \item la loi tronquée est elle-même rejetée par le bootstrap pour \code{sejour}
  ($p_{\text{tronq}} = 0{,}000$) : aucun des cinq modèles n'est adéquat pour cette
  variable ;
  \item $\alpha_{\text{séjour}} \approx 3{,}19$ est hors du domaine
  $]2, 3[$ de la relation ON/OFF, donc $H_{\text{pred}}$ n'est pas calculable
  pour la variable principale.
\end{itemize}

% =====================================================================
\section{Tableau récapitulatif des constantes}
% =====================================================================

\begin{table}[h]
\centering
\caption{Constantes du pipeline Hurst (\code{config.py}, V5)}
\label{tab:constantes}
\small
\begin{tabular}{@{}p{5cm}p{2cm}p{7cm}@{}}
\toprule
Constante & Valeur & Rôle \\
\midrule
\code{TIME_GAP_THRESHOLD} & 900~s & coupure de segment \\
\code{MIN_POINTS} & 300 & minimum de points par segment \\
\code{DISTANCE_MIN} & 0{,}1~km & seuil de distance (plancher) \\
\code{DISTANCE_MIN_RELATIF} & True & active le seuil relatif \\
\code{DISTANCE_MIN_V_PLANCHER_KMH} & 1~km/h & vitesse de marche lente \\
\code{MAX_DIST_M} & 30 & multiplicateur de $\tilde d$ \\
\code{MAX_DIST_M_ABSOLU} & 500~m & plafond absolu \\
\code{MAX_SPEED} & 200~km/h & plafond absolu de vitesse \\
\code{MAX_SPEED_RELATIF} & 20 & multiplicateur de $\tilde v$ \\
\code{MAD_Z_SEUIL} & 3{,}5 & seuil du Z-score modifié \\
\code{MAD_PERSISTANCE_MIN} & 5 & plancher de persistance \\
\code{MAD_PERSISTANCE_FRACTION} & 0{,}01 & fraction de $N$ \\
\code{MAD_DELTA_V_BRUTAL} & 20~km/h & plancher de variation brutale \\
\code{MAD_DELTA_V_FACTEUR_MAD} & 5 & multiplicateur de $\mathrm{MAD}_v$ \\
\code{MAD_FALLBACK_SPEED} & 5~km/h & plancher si MAD=0 \\
\code{MAD_FALLBACK_SPEED_FACTEUR} & 3 & multiplicateur si MAD=0 \\
\code{N_S} & 20 & nombre d'échelles DFA \\
\code{DFA_ECHELLE_MIN} & 4 & plus petite échelle \\
\code{DFA_DIVISEUR_ECHELLE_MAX} & 4 & $s_{\max} = N/4$ \\
\code{DFA_ORDRE} & 1 & ordre du détrentage \\
\code{N_M} & 20 & nombre de tailles de bloc R/S \\
\code{M_MIN_RS} & 10 & plus petite taille de bloc \\
\code{RS_DIVISEUR_M_MAX} & 2 & $m_{\max} = N/2$ \\
\code{K_WHITTLE} & 1000 & taille de la grille \\
\code{WHITTLE_EXPOSANT_M} & 0{,}65 & $m = N^{0{,}65}$ \\
\code{WHITTLE_DIVISEUR_M_MAX} & 4 & $m \leq N/4$ \\
\code{NUL_NB_SURROGATS} & 99 & surrogats par segment \\
\code{NUL_ORDRE_AR_MAX} & 5 & ordre maximal de l'AR \\
\code{NUL_ALPHA} & 0{,}05 & seuil du test nul \\
\code{SEUIL_DIVERGENCE} & 0{,}15 & seuil d'accord inter-méthodes \\
\code{SEUIL_FORT} & 0{,}5 & seuil du scénario forts \\
\bottomrule
\end{tabular}
\end{table}

% =====================================================================
\begin{thebibliography}{99}
% =====================================================================

\bibitem{peng1994} C.-K. Peng, S.~V. Buldyrev, S. Havlin, M. Simons, H.~E. Stanley,
A.~L. Goldberger. Mosaic organization of DNA nucleotides. \textit{Physical Review E},
49(2):1685--1689, 1994.

\bibitem{hurst1951} H.~E. Hurst. Long-term storage capacity of reservoirs.
\textit{Transactions of the American Society of Civil Engineers}, 116:770--799, 1951.

\bibitem{mandelbrot1968} B.~B. Mandelbrot, J.~W. Van Ness. Fractional Brownian
motions, fractional noises and applications. \textit{SIAM Review}, 10(4):422--437,
1968.

\bibitem{mandelbrot1969} B.~B. Mandelbrot, J.~R. Wallis. Robustness of the rescaled
range $R/S$ in the measurement of noncyclic long run statistical dependence.
\textit{Water Resources Research}, 5(5):967--988, 1969.

\bibitem{anis1976} A.~A. Anis, E.~H. Lloyd. The expected value of the adjusted
rescaled Hurst range of independent normal summands. \textit{Biometrika},
63(1):111--116, 1976.

\bibitem{robinson1995} P.~M. Robinson. Gaussian semiparametric estimation of long
range dependence. \textit{The Annals of Statistics}, 23(5):1630--1661, 1995.

\bibitem{phillips2004} P.~C.~B. Phillips, K. Shimotsu. Local Whittle estimation in
nonstationary and unit root cases. \textit{The Annals of Statistics}, 32(2), 2004.

\bibitem{shimotsu2005} K. Shimotsu, P.~C.~B. Phillips. Exact local Whittle
estimation of fractional integration. \textit{The Annals of Statistics}, 33(4), 2005.

\bibitem{chen2002} Z. Chen, P.~Ch. Ivanov, K. Hu, H.~E. Stanley. Effect of
nonstationarities on detrended fluctuation analysis. \textit{Physical Review E},
65:041107, 2002.

\bibitem{maraun2004} D. Maraun, H.~W. Rust, J. Timmer. Tempting long-memory --- on
the interpretation of DFA results. \textit{Nonlinear Processes in Geophysics},
11(4):495--503, 2004.

\bibitem{iglewicz1993} B. Iglewicz, D.~C. Hoaglin. \textit{How to Detect and Handle
Outliers}. ASQC Quality Press, Milwaukee, 1993.

\bibitem{leys2013} C. Leys, C. Ley, O. Klein, P. Bernard, L. Licata. Detecting
outliers: do not use standard deviation around the mean, use absolute deviation
around the median. \textit{Journal of Experimental Social Psychology},
49(4):764--766, 2013.

\bibitem{diebold2001} F.~X. Diebold, A. Inoue. Long memory and regime switching.
\textit{Journal of Econometrics}, 105(1):131--159, 2001.

\bibitem{clauset2009} A. Clauset, C.~R. Shalizi, M.~E.~J. Newman. Power-law
distributions in empirical data. \textit{SIAM Review}, 51(4):661--703, 2009.

\bibitem{vuong1989} Q.~H. Vuong. Likelihood ratio tests for model selection and
non-nested hypotheses. \textit{Econometrica}, 57(2):307--333, 1989.

\bibitem{kivela2015} M. Kivelä, M.~A. Porter. Estimating interevent time
distributions from finite observation periods in communication networks.
\textit{Physical Review E}, 92:052813, 2015.

\bibitem{taqqu1997} M.~S. Taqqu, W. Willinger, R. Sherman. Proof of a fundamental
result in self-similar traffic modeling. \textit{ACM SIGCOMM Computer
Communication Review}, 27(2):5--23, 1997.

\end{thebibliography}

\end{document}